\documentclass[11pt]{article}
\usepackage[a4paper, top=18mm, bottom=18mm, left=20mm, right=20mm]{geometry}
\usepackage[T2A]{fontenc}
\usepackage[utf8]{inputenc}
\usepackage[ukrainian]{babel}
\usepackage{graphicx}
\setlength{\parindent}{0pt}
\setlength{\parskip}{4pt}
\pagestyle{empty}
\begin{document}
%begin
{\large\textbf{Контрольна робота}}

\textbf{Варіант \No\,1}\hfill \textit{ПІБ:}\,\underline{\hspace{60mm}}
\vspace{4mm}

%beginitem
1. Що буде виведено за виконання фрагменту коду?
%code
#include <iostream>
using namespace std;

int f(int &x, int &y){
    x = 4;
    y+= 3;
    return x;
}

int main(){
    int a = 5, b = 7;
    a = f(a, a);
    cout << a << ":" << b <<':';
    {
        int a = 2, b = 6;
        cout << ((a>6) && ((b+=1) > 6))
             << a << ":" << b << ":";
    }
    {
        inta = 9;
        cout << a << ':';
    }
    cout << a << ":" << b << endl;
    return 0;
}
%code


%enditem
%beginitem
2. Що буде виведено за виконання фрагменту коду?
%code
#include <iostream>
using namespace std;

int f(int &x, int &y){
    x = 4;
    y+= 3;
    return x;
}

int main(){
    int a = 5, b = 7;
    a = f(a, a);
    cout << a << ":" << b <<':';
    {
        int a = 2, b = 6;
        cout << ((a>6) && ((b+=1) > 6)) << ':';
        cout << a << ":" << b << ':';
    }
    {
        inta = 9;
        cout << a << ':';
    }
    cout << a << ":" << b << endl;
    return 0;
}
%code


%enditem
%beginitem
3. Що буде виведено за виконання фрагменту коду?

%code
cout << (6 * 6 * 6 * 6);
%code

%enditem
%beginitem
4. Що буде виведено за виконання фрагменту коду?
%code
if (2 > 2)
    cout << "c";
else
    cout << "c";
%code


%enditem
%beginitem
5. Що буде виведено за виконання фрагменту коду?

%code
cout << (!4>4.0 && !5>= true && 2<5.0);
%code

%enditem
%beginitem
6. Що буде виведено за виконання фрагменту коду?

%code
print(-2**-2+1)
%code

%enditem
%beginitem
7. Що буде виведено за виконання фрагменту коду?

%code
#include <iostream>
using namespace std;

int f(int a, int b){
    int c = 72;
    if (a > -2)
        return 2;
    if (b >= 4)
         c = 9;
    else 
        c = 0;
    return c;
}

int main(){
    cout << f(-6, -6);
    return 0;
}
%code

%enditem
%beginitem
8. Що буде виведено за виконання фрагменту коду?

%code
#include <iostream>

int f(int c){
    int w = 72;
    if (c > -2) 
        return 2;
    if (c >= 4)
         w = 9;
    else
         w = 0;
    return w;
}

int main(){
    std::cout << f(-6);
    return 0;
}
%code


%enditem
%beginitem
9. Що буде виведено за виконання фрагменту коду?

%code
cout << 6 * 6 * 6 * 6;
%code

%enditem
%beginitem
10. Відмітити всі правильні варіанти:
( 1 ) 8,3

( 2 ) <>

( 3 ) 8.3

( 4 ) FIFO

%enditem
%beginitem
11. Що буде виведено за виконання фрагменту коду?

%code
cout << (4 > 4.0 >= 5 < true);
%code

%enditem
%beginitem
12. Що буде виведено за виконання фрагменту коду?

%code
cout << (6 * 6 * 6 * 6);
%code

%enditem
%beginitem
13. Що буде виведено за виконання фрагменту коду?

%code
#include <iostream>
using namespace std;

int f(int c){
    int w = 72;
    if (c > -2) 
        return 2;
    if (c >= 4)
         w = 9;
    else
         w = 0;
    return w;
}

int main(){
    cout << f(-6);
    return 0;
}
%code


%enditem
%beginitem
14. Відмітити все, що є коректним оператором С++:
( 1 ) cout>>3;

( 2 ) int x+y=z;

( 3 ) cout<<3.4;

( 4 ) int f(int n);

%enditem
%beginitem
15. Що буде виведено за виконання фрагменту коду?
%code
if (2 > 2)
    cout << "c";
else
    cout << "v";
%code


%enditem
%beginitem
16. Що буде виведено за виконання фрагменту коду?

%code
#include <iostream>
using namespace std;

int a = 2, b = 9, c = 0;

int f(int &a){
int c;    a *= 4;
    b = 3;
    c = 5;
    return a + b + c;
}

int main(){
    inta = 5;
    intb = 4;
    intc = 6;
    cout << f(a) << ':';
    cout << a << ':' << b << ':' << c;
    return 0; 
}
%code

%enditem
%beginitem
17. Що буде виведено за виконання фрагменту коду?

%code
cout << 6 * 6 * 6 * 6;
%code

%enditem
%beginitem
18. Що буде виведено за виконання фрагменту коду?

%code
cout << (!4 > 4.0 && !5 >= true && 2 < 5.0);
%code

%enditem
%beginitem
19. Що буде виведено за виконання фрагменту коду?

%code
print(-2**-2+1)
%code

%enditem
%beginitem
20. Що буде виведено за виконання фрагменту коду?

%code
#include <iostream>
using namespace std;

int a = 2, b = 9, c = 0;

int f(int &a){
int c;    a *= 4;
    b = 3;
    c = 5;
    return a + b + c;
}

int main(){
    inta = 5;
    intb = 4;
    intc = 6;
    cout << f(a) << ':';
    cout << a << ':' << b << ':' << c;
    return 0; 
}
%code

%enditem
%beginitem
21. Що буде виведено за виконання фрагменту коду?

%code
#include <iostream>
using namespace std;

bool f(int n){
    cout<<"f";
    return n;
}

int main(){
    cout<<(f(-6) && f(0));
    return 0;
}
%code

%enditem
%beginitem
22. Відмітити всі правильні варіанти:
( 1 ) 8,3

( 2 ) <>

( 3 ) 8.3

( 4 ) FIFO

%enditem
%beginitem
23. Що буде виведено за виконання фрагменту коду?

%code
cout << (4 > 4.0 >= 5 < true);
%code

%enditem
%beginitem
24. Відмітити все, що є коректним оператором С++:
( 1 ) cout>>3;

( 2 ) int x+y=z;

( 3 ) cout<<3.4;

( 4 ) int f(int n);

%enditem
%beginitem
25. Що буде виведено за виконання фрагменту коду?

%code
#include <iostream>

int f(int a, int b){
    int c = 72;
    if (a > -2)
        return 2;
    if (b >= 4)
         c = 9;
    else 
        c = 0;
    return c;
}

int main(){
    std::cout << f(-6, -6);
    return 0;
}
%code

%enditem
%beginitem
26. Що буде виведено за виконання фрагменту коду?

%code
#include <iostream>
using namespace std;

int a = 2, b = 9, c = 0;

int f(){
    inta = 2;
    intb = 5;
    intc = 4;
    return a + b + c;
}

int main(){
    inta = 6;
    intb = 7;
    intc = 1;
    cout << f() << ':';
    cout << a << ':' << b << ':' << c;
    return 0;
}
%code

%enditem
%beginitem
27. Що буде виведено за виконання фрагменту коду?

%code
#include <iostream>

bool f(int n){
    std::cout<<"f";
    return n;
}

int main(){
    std::cout<<(f(-6) && f(0));
    return 0;
}
%code

%enditem
%beginitem
28. Що буде виведено за виконання фрагменту коду?

%code
#include <iostream>
using namespace std;

int a = 2, b = 9, c = 0;

int f(){
    inta = 2;
    intb = 5;
    intc = 4;
    return a + b + c;
}

int main(){
    inta = 6;
    intb = 7;
    intc = 1;
    cout << f() << ':';
    cout << a << ':' << b << ':' << c;
    return 0;
}
%code

%enditem




























%end
\newpage
\end{document}

